\section*{Capítulo \ref{ch:g9:weight-vs-mass} --- Peso frente a masa: medir fuerzas}

\begin{solution}{exo:g9:weight-vs-mass:1}
El \omterm{def:g9:weight-vs-mass:newton}{newton} (\unit{N}), medido con el \omterm{def:g9:weight-vs-mass:dynamometer}{dinamómetro}. Anclas:
un \omterm{def:g9:weight-vs-mass:newton}{newton} --- el \omterm{def:g4:weight-and-mass:weight}{peso} de una manzana pequeña; cien ---
un apretón de manos firme (o el \omterm{def:g4:weight-and-mass:weight}{peso} de una maleta de
\qty{10}{kg}).
\end{solution}

\begin{solution}{exo:g9:weight-vs-mass:2}
$P = m \times g$: el \omterm{def:g4:weight-and-mass:weight}{peso} en \omterm{def:g9:weight-vs-mass:newton}{newtons}, la
\omterm{def:g3:measuring-mass:mass}{masa} en \omterm{def:g3:measuring-mass:units}{kilogramos} y $g$ --- la
\omterm{prop:g9:weight-vs-mass:formula}{intensidad gravitatoria} del lugar --- los \omterm{def:g9:weight-vs-mass:newton}{newtons} de tirón
que recibe allí cada \omterm{def:g3:measuring-mass:units}{kilogramo} (\qty{9.8}{N/kg} en la Tierra).
\end{solution}

\begin{solution}{exo:g9:weight-vs-mass:3}
Perro: $25 \times 9.8 = \qty{245}{N}$. Hogaza: $0.7 \times 9.8 \approx
\qty{6.9}{N}$. Un lector de \qty{55}{kg}: unos \qty{539}{N}.
\end{solution}

\begin{solution}{exo:g9:weight-vs-mass:4}
$m = 14.7 \div 9.8 = \qty{1.5}{kg}$.
\end{solution}

\begin{solution}{exo:g9:weight-vs-mass:5}
La pendiente es el $g$ local --- \omterm{def:g9:weight-vs-mass:newton}{newtons} por
\omterm{def:g3:measuring-mass:units}{kilogramo}. En la \omterm{def:g2:moon-in-the-sky:moon}{Luna}, la recta sigue siendo recta y
pasando por el origen, pero se desploma hasta una pendiente suave de
$1.6$: la misma proporcionalidad, un mundo más flojo.
\end{solution}

\begin{solution}{exo:g9:weight-vs-mass:6}
La columna de la \omterm{def:g3:measuring-mass:mass}{masa} no cambia nunca --- \qty{10}{kg} de materia
viajan intactos. Las columnas de $g$ y del \omterm{def:g4:weight-and-mass:weight}{peso} cambian con el
mundo: el \omterm{def:g4:weight-and-mass:weight}{peso} es el peaje local, $m \times g$.
\end{solution}

\begin{solution}{exo:g9:weight-vs-mass:7}
La \omterm{def:g3:levers-and-scales:scale}{balanza} de platillos compara $m_1 g$ con $m_2 g$: $g$ multiplica los
dos platillos y se cancela --- \omterm{def:g3:measuring-mass:mass}{masa}, con verdad, en cualquier
sitio. La báscula de baño mide $P = m g_{\text{aquí}}$ pero divide
entre $g_{\text{Tierra}}$ para imprimir \omterm{def:g3:measuring-mass:units}{kilogramos}: en la
\omterm{def:g2:moon-in-the-sky:moon}{Luna} imprime $m \times 1.6/9.8$ --- una sexta parte de la verdad.
\end{solution}

\begin{solution}{exo:g9:weight-vs-mass:8}
\omterm{def:g4:weight-and-mass:weight}{Peso}: \qty{19.6}{N} hacia abajo --- una flecha de unos
\qty{3.9}{cm} a cinco \omterm{def:g9:weight-vs-mass:newton}{newtons} por \omterm{def:g2:measuring-length:units}{centímetro}, desde el
libro hacia abajo. Sostén: una flecha igual de \qty{3.9}{cm} hacia
arriba. Longitudes iguales, sentidos opuestos: presupuesto cuadrado,
libro quieto.
\end{solution}

\begin{solution}{exo:g9:weight-vs-mass:9}
El avión tiene que llevar, levantar y frenar la \emph{materia} del
bulto --- su \omterm{def:g3:measuring-mass:mass}{masa} --- y \qty{23}{kg} de materia son
\qty{23}{kg} en cualquier mundo. El mostrador lunar necesita una
\omterm{def:g3:levers-and-scales:scale}{balanza} de platillos con \omterm{def:g3:measuring-mass:mass}{masas} marcadas: la única que sobrevive al
cambio de $g$.
\end{solution}

\begin{solution}{exo:g9:weight-vs-mass:10}
$m = 111 \div 3.7 = \qty{30}{kg}$; en la Tierra pesa $30 \times 9.8 =
\qty{294}{N}$.
\end{solution}

\begin{solution}{exo:g9:weight-vs-mass:11}
Marca de más: colgando, el muelle carga también con su propio
\omterm{def:g4:weight-and-mass:weight}{peso} (y con el de su gancho), que el cero tumbado no incluía nunca
--- el instrumento añade un excedente constante. Ponlo a cero colgando
y el excedente queda restado antes de que llegue la carga.
\end{solution}

\begin{solution}{exo:g9:weight-vs-mass:12}
Material: \omterm{def:g9:weight-vs-mass:dynamometer}{dinamómetro}, juego de \omterm{def:g3:measuring-mass:mass}{masas} marcadas (de $1$ a
\qty{5}{kg}) y gancho. Cuelga cada \omterm{def:g3:measuring-mass:mass}{masa}, anota las parejas ($m$,
$P$) y representa $P$ frente a $m$: los puntos se alinean pasando por
el origen, y la pendiente --- lo que sube en \omterm{def:g9:weight-vs-mass:newton}{newtons} entre lo que
avanza en \omterm{def:g3:measuring-mass:units}{kilogramos} --- es $g \approx \qty{9.8}{N/kg}$. En la
\omterm{def:g2:moon-in-the-sky:moon}{Luna}: protocolo idéntico, rectitud idéntica y pendiente nueva de
$1.6$ --- el método mide el mundo en el que esté; lo único local
es el número.
\end{solution}

\begin{solution}{pb:g9:weight-vs-mass:1}
\textbf{1.} Tierra $70 \times 9.8 = \qty{686}{N}$; \omterm{def:g2:moon-in-the-sky:moon}{Luna}
\qty{112}{N}; Marte \qty{259}{N}; cubierta de Júpiter
\qty{1736}{N}.
\textbf{2.} Correcta en la Tierra: \qty{70}{kg}. \omterm{def:g2:moon-in-the-sky:moon}{Luna}: $112 \div
9.8 \approx \qty{11.4}{kg}$; Marte: $259 \div 9.8 \approx
\qty{26.4}{kg}$; Júpiter: $1736 \div 9.8 \approx \qty{177}{kg}$ ---
adulación y difamación por fórmula.
\textbf{3.} \qty{70}{kg} en todos los puertos: la comparación de la
\omterm{def:g3:levers-and-scales:scale}{balanza} cancela el $g$ local.
\textbf{4.} Solo la cubierta de Júpiter: los \qty{1736}{N} se acercan
al límite pero se quedan por debajo --- ningún puerto rompe los
\qty{2000}{N} para este viajero (un compañero más pesado, de
\qty{90}{kg}, sí lo rompería allí: $90 \times 24.8 =
\qty{2232}{N}$).
\textbf{5.} Salida: $40 \times 9.8 = \qty{392}{N}$; llegada: $40
\times 3.7 = \qty{148}{N}$.
\textbf{6.} $m = 200/g$: Tierra \qty{20.4}{kg}; \omterm{def:g2:moon-in-the-sky:moon}{Luna}
\qty{125}{kg}; Marte \qty{54}{kg}; Júpiter \qty{8.1}{kg} --- los
\omterm{def:g9:weight-vs-mass:newton}{newtons} fijos del sindicato compran \omterm{def:g3:measuring-mass:mass}{masas} muy distintas.
\textbf{7.} $m = 150 \div 1.6 = \qty{93.75}{kg}$ --- unos
\qty{94}{kg} de caja para la aduana.
\textbf{8.} $m = 294 \div 9.8 = \qty{30}{kg}$; \omterm{def:g4:weight-and-mass:weight}{peso} en la
\omterm{def:g2:moon-in-the-sky:moon}{Luna} $30 \times 1.6 = \qty{48}{N}$.
\textbf{9.} Los músculos y la \omterm{def:g3:measuring-mass:mass}{masa} volvieron a casa sin cambios;
lo único que cambió fue el adversario. En la \omterm{def:g2:moon-in-the-sky:moon}{Luna}, cada
\omterm{def:g3:measuring-mass:units}{kilogramo} costaba \qty{1.6}{N} de levantar; en casa vuelven a
cobrarle \qty{9.8}{N}. Las ganancias eran el descuento de la
\omterm{def:g2:moon-in-the-sky:moon}{Luna}, no el crecimiento del cuerpo --- $P = m g$, con $m$ leal y
$g$ voluble.
\textbf{10.} $0.2 \times 24.8 \approx \qty{5}{N}$ --- un yunque no es:
en la Tierra las varillas pesaban \qty{2}{N}, así que batir en la
cubierta es dos veces y media más pesado, cansado para una muñeca pero
poco parecido al trabajo de una fundición. Queja: estimada, en parte.
\textbf{11.} Un \omterm{def:g9:weight-vs-mass:newton}{newton} de chocolate es $m = 1/g$: en la Tierra,
unos \qty{102}{g}; en la \omterm{def:g2:moon-in-the-sky:moon}{Luna}, \qty{625}{g}; en Júpiter,
\qty{40}{g}. Compra tus \omterm{def:g9:weight-vs-mass:newton}{newtons} en la \omterm{def:g2:moon-in-the-sky:moon}{Luna}.
\textbf{12.} Por ejemplo: ``La \omterm{def:g3:measuring-mass:mass}{masa} viaja contigo; el
\omterm{def:g4:weight-and-mass:weight}{peso} es el peaje que le cobra cada mundo --- \omterm{def:g6:light-propagation:ray}{haz} la maleta en
\omterm{def:g3:measuring-mass:units}{kilogramos} y deja que los puertos cobren sus
\omterm{def:g9:weight-vs-mass:newton}{newtons}''.
\end{solution}
